\section{Zadanie 9}
Mając do dyspozycji generator liczb losowych z rozkłady jednostajnego na odcinku \([0,1]\)
zaprojektuj generator liczb z rozkładu Cauchy'ego

\subsection*{Rozwiązanie}
W tym zadaniu skorzystamy z metody odwrotnej dystrybuanty.\\
Standardowy rozkład Cauchy'ego ma funkcję gęstości prawdopodobieństwa:
\[ f(x) = \frac{1}{\pi(1+x^2)} \]
Znajdujemy dystrybuantę:
\[ F(x) = \int_{-\infty}^{x}f(t)dt = \int_{-\infty}^{x}\frac{1}{\pi(1+t^2)}dt \]
\[ F(x) = \frac{1}{\pi}\arctan x + \frac{1}{2}, \qquad x \in \mathbb{R} \]
Jest to funkcja ściśle rosnąca, a zatem posiada funkcję odwrotną:
\[ F^{-1}(u) = \tan\!\left(\pi(u - \tfrac{1}{2})\right), \qquad u \in [0,1] \]

Niech \(U\) będzie zmienną losową o rozkładzie jednostajnym na przedziale \([0,1]\).
Definiujemy zmienną losową
\[
	X = F^{-1}(U) = \tan\!\left(\pi(U - \tfrac{1}{2})\right).
\]
Sprawdźmy czy genertor działa, rzeczywiście dla dowolnego \(x \in \mathbb{R}\) zachodzi
\[
	\begin{aligned}
		\mathbb{P}(X \le x)
		 & = \mathbb{P}\!\left( \tan\!\left(\pi(U - \tfrac{1}{2})\right) \le x \right) \\
		 & = \mathbb{P}\!\left( U \le \frac{1}{\pi}\arctan x + \frac{1}{2} \right)     \\
		 & = \frac{1}{\pi}\arctan x + \frac{1}{2}
	\end{aligned}
\]
co oznacza, że zmienna losowa \(X\) ma standardowy rozkład Cauchy'ego. \\



\subsection*{Odpowiedź}
\[
	\boxed{X = \tan\!\left(\pi(U - \tfrac{1}{2})\right).}
\]
